% Example: Numbered equations with \begin{equation}
% Demonstrates how the equation environment assigns a reference number
% to a displayed formula, and how \label / \eqref can cross-reference it.
\documentclass[12pt, a4paper]{article}
\usepackage{amsmath, amssymb, amsthm}
\usepackage{fontspec}
\begin{document}

\section*{Numbered Equations}

The \texttt{equation} environment numbers each formula automatically:

\begin{equation}
  E = mc^2
  \label{eq:einstein}
\end{equation}

\begin{equation}
  \nabla \times \mathbf{E} = -\frac{\partial \mathbf{B}}{\partial t}
  \label{eq:faraday}
\end{equation}

\begin{equation}
  i\hbar \frac{\partial}{\partial t}\Psi(\mathbf{r},t) =
  \hat{H}\Psi(\mathbf{r},t)
  \label{eq:schrodinger}
\end{equation}

We can refer to these equations by number. Equation~\eqref{eq:einstein}
expresses mass--energy equivalence. Equation~\eqref{eq:faraday} is
Faraday's law of induction, and Equation~\eqref{eq:schrodinger} is the
time-dependent Schr\"odinger equation.

To suppress the number for a single equation, use \texttt{equation*}:

\begin{equation*}
  \frac{d}{dx}\int_a^x f(t)\,dt = f(x)
\end{equation*}

\end{document}
